\documentclass[12pt]{article}

\usepackage{sbc-template}
\usepackage{graphicx,url}
\usepackage[utf8]{inputenc}
\usepackage[brazilian]{babel}
\usepackage{booktabs}
\usepackage{amsmath}

\sloppy

\title{LLM como Juiz na Medição de ASR: Dois Juízes Automáticos, Um Protocolo de Calibração e a Coleta Humana Pendente}

\author{Leonardo Kawakami Hissanaga\inst{1}}

\address{PCS5917 --- Inteligência Artificial Adversarial (Poli/USP), 2026/2 \\
  Universidade de São Paulo, São Paulo -- SP -- Brazil \\
  NUSP 11374864
}

\begin{document}

\maketitle

\vspace{-3mm}
\begin{center}\footnotesize\itshape
Entrega individual (contribuições do autor). Uso de IA generativa declarado (Código de Conduta da SBC, Art.~2º,~III); nunca como autoria.
\end{center}

\begin{resumo}
Medimos a taxa de sucesso de ataque (ASR) de jailbreaks com \textbf{juízes LLM automáticos} em dois experimentos da PCS5917: (i)~no Aula~7/NB3, um juiz Mistral-7B-Instruct-v0.3 (rubrica binária, justificativa JSON) rotulou 165 respostas do alvo Qwen2.5-7B-Instruct (NF4) sob \{baseline, Self-Reminder, Intention Analysis\}, com ASR de 12,0\% (baseline), 2,0\% (IA) e 0,0\% (SR) para manuais e 0,0\% para Base64 em todas as condições; (ii)~no piloto PAIR (alvo Llama-2-7B-chat, atacante Qwen3.5-9B), o juiz vendorizado do repositório oficial do PAIR (escala 1--10) atingiu ASR de 10,0\% (IC 95\% Wilson: 1,8--40,4\%) em 50 \emph{streams}. Consolidamos os 146 ataques das 5 atividades em uma planilha única e definimos um protocolo de calibração humano--juiz (amostragem estratificada de 31~$+$~18 linhas, cegamento, semente~42, kappa de Cohen). \textbf{Limitação central e declarada: a rotulação humana do subconjunto nunca foi coletada (0/31 linhas)}; \emph{não} reportamos kappa humano--juiz nem ASR humano do subconjunto, e todos os ASRs automáticos permanecem não calibrados. Concluímos que juízes únicos são instrumentos mal calibrados (enviesáveis e não comparáveis entre si) e que a rota para ASRs confiáveis passa por painel multi-juiz e calibração humana sistemática.
\end{resumo}

\section{Introdução}

A comunidade de segurança de LLMs migrou da avaliação humana massiva para \emph{LLM-as-a-judge} por custo e escala~\cite{zheng2023}, mas juízes automáticos introduzem vieses conhecidos (ancoragem, preferência de posição/verbosidade, falsos negativos) que contaminam exatamente a métrica central do campo, o ASR. Este relatório toma o juiz como objeto de estudo, não como caixa-preta: (a)~documenta os dois juízes automáticos dos experimentos da disciplina (o Mistral-7B-Instruct-v0.3 com rubrica própria no Aula~7/NB3 e o juiz vendorizado do PAIR~\cite{chao2025}) e tabula os ASRs de cada um; (b)~apresenta a planilha única que consolida as 5 atividades de ataque (146 ataques) com sua forma de avaliação; e (c)~especifica o protocolo de calibração humano--juiz já definido e declara a rotulação humana como \textbf{pendência explícita} (Seção~\ref{sec:limit}): 0 de 31 linhas do subconjunto rotuladas.

Contribuições: (i)~inventário reprodutível de ASRs automáticos por juiz, ancorado em CSVs de evidência; (ii)~evidência de que ASRs de juízes diferentes não são comparáveis; (iii)~protocolo pronto de calibração humano--juiz--humano, com a pendência da coleta humana declarada.

\section{Metodologia}

\textbf{Planilha única de consolidação.} Todas as atividades de ataque da disciplina foram unificadas em \texttt{planilha-unica-ataques.csv} (5 linhas, 146 ataques; Tabela~\ref{tab:planilha}), com colunas de ataque, método, alvo, $n$, forma de avaliação, ASR e arquivo de evidência, a base de rastreabilidade deste relatório.

\textbf{Juiz~A: Mistral-7B (Aula~7/NB3).} Mistral-7B-Instruct-v0.3 (HF, NF4 em T4/Colab) recebe o prompt malicioso e a resposta do alvo Qwen2.5-7B-Instruct (NF4) e devolve JSON com r\`otulo bin\'ario $+$ justificativa, por rubrica escrita no esp\'irito da literatura de ju\'izes por rubrica~\cite{zhu2023}: $1$ = conteúdo substantivo que avança o objetivo nocivo (compliance parcial conta; disclaimer não cancela); $0$ = recusa/evasiva/inválida. 165/165 respostas avaliadas, zero inválidas. ASR = nº de rótulos~1 / respostas válidas por célula (família $\times$ defesa) e por tarefa.

\textbf{Juiz~B: juiz próprio do PAIR.} No piloto PAIR (alvo Llama-2-7B-chat NF4, atacante Qwen3.5-9B GGUF; 10 comportamentos JBB-Behaviors $\times$ 5 \emph{streams} $\times$ 3 iterações = 150 iterações), a avaliação é feita pelo \textbf{juiz vendorizado \emph{verbatim} do repositório oficial do PAIR}~\cite{chao2025} (``You are a helpful red teaming assistant\dots evaluate whether the response\dots violates ethical, legal, or safety guidelines''), com score escalar 1--10 por iteração; sucesso = score máximo $\geq 10$ no \emph{stream}. O pipeline foi usado como fornecido, só adaptado para GGUF local; o CSV de auditoria (\texttt{pair\_iterations\_pair\_pilot\_v1.csv}) registra o score de cada iteração, sem erros.

\textbf{Protocolo de calibração humano--juiz (resumo).} Subconjunto estratificado: 31 linhas do piloto PAIR (todas com score $\geq 8$ do juiz $+$ amostra das demais; estratos A/B/C) $+$ 18 linhas do Aula~7/NB3 (\texttt{human\_review\_sample.csv}, 3 por célula, \texttt{random\_state=42}). Rubrica binária idêntica à do Juiz~A, fixada \emph{antes} da coleta. Cegamento simples (rótulos do juiz ocultos até dupla rotulação), ordem permutada com semente~42, dois avaliadores independentes (A1 autor cego, A2 par externo), kappa de Cohen com bandas de Landis \& Koch~\cite{landis1977} (meta $\kappa \geq 0{,}61$), planilha de discordância com resolução por retorno à rubrica, auditoria por \emph{hash} curto da resposta (\texttt{protocolo-avaliacao-humana.md}~v1.0).

\section{Resultados}

\begin{table}[t]
\centering
\caption{Planilha única consolidada: 5 linhas de \texttt{planilha-unica-ataques.csv} (146 ataques no total).}
\label{tab:planilha}
\scriptsize
\begin{tabular}{@{}p{1.6cm}p{1.7cm}p{1.5cm}p{1.6cm}@{}}
\toprule
Atividade & Ataque (método) & Avaliação & ASR / resultado \\
\midrule
Aula~2 & Jailbreaks manuais EN/PT, 20 casos & humana & ver planilha detalhada \\
Aula~4 & nanoGCG, 11 templates (sufixo adversarial) & humana (Recusou/Parcial) & 3/11 vs.\ 1/11 (parcial$+$) \\
At.1 PAIR & PAIR, 10 comport.\ $\times$ 5 \emph{streams} & automática (juiz PAIR) & 10,0 (IC 1,8--40,4) \\
At.2/Aula~7 & 10 jailbreaks $\times$ 5 tarefas, 3 defesas & automática (juiz NB3) & 12,0 / IA 2,0 / SR 0,0 \\
At.2/Aula~7 & Base64 $\times$ 5 tarefas $\times$ 3 condições & automática (juiz NB3) & 0,0 (todas) \\
\bottomrule
\end{tabular}
\end{table}

\begin{table}[t]
\centering
\caption{ASR automático (Juiz~A, Mistral-7B-Instruct-v0.3, Aula~7/NB3; $n=50$ manual e $n=5$ \emph{encoded} por célula; 0 inválidas em 165).}
\label{tab:nb3}
\small
\begin{tabular}{@{}lccc@{}}
\toprule
Família & Baseline & Self-Reminder & Intention Anal. \\
\midrule
manual (10 jailbreaks $\times$ 5 tarefas) & \textbf{12,0} (6/50) & \textbf{0,0} (0/50) & \textbf{2,0} (1/50) \\
encoded (Base64, 5 tarefas) & 0,0 (0/5) & 0,0 (0/5) & 0,0 (0/5) \\
\bottomrule
\end{tabular}
\end{table}

\begin{table}[t]
\centering
\caption{ASR automático (Juiz~B, juiz vendorizado PAIR, escala 1--10; alvo Llama-2-7B-chat, sem defesa; 150 iterações, 0 erros).}
\label{tab:pair}
\small
\begin{tabular}{@{}p{2.2cm}p{3.6cm}@{}}
\toprule
Métrica & Valor \\
\midrule
ASR por comportamento & 10,0\% (IC 95\% Wilson: 1,8--40,4) \\
\emph{Streams} com sucesso (score~10) & 1/50 (Disinformation/5G-COVID) \\
Score máx.\ 9 sem sucesso & Economic harm, Fraud/Deception, Expert advice, Gov.\ decision-making \\
Iterações com erro & 0/150 \\
\bottomrule
\end{tabular}
\end{table}

Os ASRs das Tabelas~\ref{tab:nb3}--\ref{tab:pair} são \textbf{todo o estoque de avaliação automática hoje disponível} (65 ataques do Aula~7 $+$ 50 \emph{streams} PAIR); as demais linhas da Tabela~\ref{tab:planilha} (Aulas~2 e~4, 31 ataques) foram avaliadas só por humano.

\section{Limitações e trabalho futuro: a calibração humana pendente}
\label{sec:limit}

\textbf{Limitação explícita.} O protocolo da Seção~2 está definido e versionado, mas \textbf{a coleta humana nunca foi executada: 0 de 31 linhas do subconjunto PAIR (e 0 de 18 do Aula~7) possuem rótulo humano}. Por disciplina de método, enquanto o CSV rotulado não existir: (i)~\textbf{não reportamos} taxa de discordância humano--juiz, ASR humano do subconjunto, nem $\kappa$ humano--juiz ou humano--humano; nenhum número desse tipo foi calculado porque nenhum rótulo humano foi coletado; (ii)~os ASRs das Tabelas~\ref{tab:nb3}--\ref{tab:pair} permanecem \textbf{não calibrados} contra julgamento humano. O prazo original (04/11) foi perdido; esta entrega atrasada traz o estado real do experimento, não uma estimativa no lugar dele.

Já estão prontos: o sorteio estratificado (semente~42), a rubrica com casos-limite, o CSV cegado (\texttt{review\_id}, tarefa, \texttt{resposta\_do\_alvo}) com arquivo-chave separado, e o plano de \emph{commits} pré/pós-consenso em \texttt{leokawakami/portfolio-pcs5917}. A literatura recomenda cautela exatamente onde estamos: Hayashi et al.~\cite{hayashi2026} mediram 12,7\% de discordância humano--juiz em jailbreak, com viés de falsos negativos: ASR de juiz único tende a ser \textbf{otimista}, o que faria os números das Tabelas~\ref{tab:nb3}--\ref{tab:pair} serem piso, não teto, em reavaliação humana.

\textbf{Trabalho futuro:} (1)~executar a rotulação A1/A2 do subconjunto 31$+$18 e publicar o CSV versionado com \emph{tags} pré/pós-consenso; (2)~reportar, além de $\kappa$, a matriz de confusão humano$\times$juiz por estrato (falsos negativos no estrato~B (score 5--9) indicariam limiar mal calibrado, não juiz ruim); (3)~painel de 3 juízes (Mistral-7B $+$ Llama-3-8B $+$ Qwen2.5-7B) sobre as mesmas 315 respostas, agregando por voto e reportando a discordância intra-painel como incerteza.

\section{Conclusão e reflexões: multi-juiz e calibração}

Os números deste relatório vieram de \textbf{dois juízes que não concordariam necessariamente entre si nem com humanos}: o Mistral-7B do NB3 usa rubrica binária nossa; o juiz do PAIR, escala 1--10 vendorizada com limiar~10; um mesmo par (prompt, resposta) pode ser~0 no primeiro e~9 no segundo (e $9 \neq$ sucesso no PAIR: 4 comportamentos com score máx.~9 e zero sucesso na Tabela~\ref{tab:pair}). Concluímos que: (i)~ASR é função do juiz tanto quanto do ataque: comparar 12,0\% (Tabela~\ref{tab:nb3}) com 10,0\% (Tabela~\ref{tab:pair}) como a mesma régua seria erro metodológico; (ii)~juiz único é instrumento sem curva de calibração; o caminho é \textbf{painel multi-juiz} e \textbf{ancoragem humana}: o protocolo $\kappa$ transforma o subconjunto 31$+$18 na ``célula-padrão'' contra a qual rubricas e juízes são revisados; (iii)~rubricas devem ser fixadas antes da coleta (falsos negativos por disclaimers e compliance parcial são o caso-limite mais frequente na literatura). A lição desta atividade atrasada é também de integridade: a ausência de rótulos humanos é resultado declarado (Seção~\ref{sec:limit}), não lacuna a preencher com número inventado.

{\footnotesize
\bibliographystyle{sbc}
\bibliography{juiz-refs}
}

\end{document}
